\documentclass[11pt]{article}
\usepackage[T1]{fontenc}
\usepackage{lmodern,microtype}
\usepackage[margin=1in]{geometry}
\usepackage{amsmath,amssymb,amsthm,mathtools,booktabs}
\usepackage[numbers,sort&compress]{natbib}
\usepackage[colorlinks=true,linkcolor=blue,citecolor=blue,urlcolor=blue]{hyperref}
\usepackage[nameinlink,noabbrev]{cleveref}
\hypersetup{pdftitle={Classical capacity and entropy inequalities for generalized amplitude damping},pdfauthor={OpenAI}}
\urlstyle{same}
\newtheorem{theorem}{Theorem}[section]
\newtheorem{lemma}[theorem]{Lemma}
\newtheorem{proposition}[theorem]{Proposition}
\newtheorem{corollary}[theorem]{Corollary}
\theoremstyle{remark}
\newtheorem{remark}[theorem]{Remark}
\DeclareMathOperator{\Tr}{Tr}
\DeclareMathOperator{\diag}{diag}
\DeclareMathOperator{\atanh}{atanh}
\newcommand{\A}{\mathcal A}
\newcommand{\ent}{\mathcal S}
\newcommand{\C}{\mathbb C}
\newcommand{\R}{\mathbb R}
\newcommand{\ket}[1]{\lvert #1\rangle}
\newcommand{\bra}[1]{\langle #1\rvert}
\newcommand{\proj}[1]{\ket{#1}\bra{#1}}
\newcommand{\hs}[1]{\lVert #1\rVert_{\mathrm{HS}}}
\newcommand{\dd}{\,\mathrm d}
\newcommand{\doi}[1]{\href{https://doi.org/#1}{\nolinkurl{#1}}}
\numberwithin{equation}{section}

\title{Classical capacity and entropy inequalities\\for generalized amplitude damping}
\author{OpenAI}
\date{September 24, 2026}

\begin{document}
\maketitle
\begin{abstract}
We determine the unassisted classical capacity of every qubit generalized
amplitude-damping channel, for all damping strengths and thermal
occupations. It equals the one-shot Holevo capacity and is given by a
one-variable maximum. A binary pure-state ensemble attains the one-use
optimum, and its independent products attain the optimum at every block
length. We also prove that one-shot Holevo capacity, minimum output entropy,
and regularized unassisted classical capacity are additive under tensor
product with every finite-dimensional completely positive trace-preserving
partner channel.
\end{abstract}

\section{Introduction}

Amplitude damping models a qubit that loses an excitation to its
surroundings. At nonzero stationary excited-state population, transitions
in both directions are possible; generalized amplitude damping describes
this relaxation together with the decay of coherence. We determine how
much classical information can be transmitted through repeated,
independent uses of this channel when a sender may encode each message
into an arbitrary state of the entire input block. We also determine
how its one-shot and regularized classical capacities behave when it is
used in parallel with an arbitrary finite-dimensional partner channel.

This permission is essential. Although a channel acts independently at
each site, an input signal can be entangled across sites, and the receiver
can measure the outputs collectively. The pure-signal coding theorem
of Hausladen, Jozsa, Schumacher, Westmoreland, and Wootters
\cite{HausladenEtAl1996} established the role of collective decoding.
The mixed-signal coding theorems of Holevo and Schumacher--Westmoreland
\cite{Holevo1998,SchumacherWestmoreland1997}, applied to arbitrary input
blocks, express the unassisted classical capacity as a regularization of the
Holevo information over arbitrarily long blocks. Optimizing a single
use gives an achievable rate, but does not by itself rule out an
advantage from entangled signals. General Holevo additivity fails
\cite{Hastings2009}, so removing regularization requires an argument
specific to the channel.

For ordinary amplitude damping, Bennett, Shor, Smolin, and Thapliyal
\cite[preprint, Section~III~B, Eq.~(89)]{BennettEtAl2002} described the one-parameter
optimization over two equiprobable pure signals with opposite phases.
Giovannetti and Fazio
\cite[Section~III~A, Eqs.~(42)--(43)]{GiovannettiFazio2005} explicitly
derived the optimized one-use expression and its attaining ensemble.
Leditzky, Kaur, Datta, and Wilde
\cite{LeditzkyEtAl2018} identified its unrestricted classical capacity
as an open problem in their 2018 study of approximate additivity.
The geometry of two-signal qubit optimizers was developed by Cortese
\cite[Section~7.1]{Cortese2002} and Berry \cite{Berry2005}.
Hou and Fang \cite{HouFang2007} characterized the generalized channel's
one-use Holevo optimum and evaluated the resulting scalar optimization
numerically. Khatri, Sharma, and Wilde
\cite[Section~VI]{KhatriEtAl2020} placed that characterization alongside
upper bounds on the unrestricted capacity. More recently, Fang
\cite[Theorem~1]{Fang2026} derived a close converse bound throughout the
generalized family; that bound leaves a gap from the one-use value in
general.

Tang, Zhu, Bai, and Wang
\cite[Corollary~5.1, Lemma~5.2, and Theorem~5.5]{TangEtAl2026}
established Holevo and classical-capacity additivity with arbitrary
finite-dimensional partners for qubit channels admitting a pure output,
including ordinary amplitude damping. Their triangular block entropy
inequality \cite[Theorem~4.2]{TangEtAl2026} is the zero-block estimate
recalled below. At positive damping and strictly interior thermal
occupation, the generalized channel has no pure output, so that channel
criterion does not apply. We extend the triangular estimate to the
thermal block structure by a convex mixture that preserves its scalar
entropy term. Both zero-block factorizations in this mixture have the
same scalar contribution; entropy concavity therefore preserves the
lower bound needed for the thermal channel.

Some parameter regimes already have exact capacity formulas from
broader additivity theorems. King's theorem for unital qubit channels
\cite{King2002} applies at stationary occupation $1/2$, and Shor's
entanglement-breaking theorem \cite{Shor2002} applies in the corresponding
region of the generalized family; see
\cite[Section~VI]{KhatriEtAl2020} for that specialization. Those theorems
permit an arbitrary partner channel. The entropy argument below covers
every parameter pair. Its rectangular form also gives additivity of
one-shot Holevo capacity and minimum output entropy with every
finite-dimensional partner, and regularization gives additivity of
unassisted classical capacity for the same pair.

\subsection{Definitions and result}

All vector spaces are finite dimensional, and a star denotes the
adjoint. Write $\mathcal B(\mathcal H)$ for the algebra of linear
operators on a complex Hilbert space $\mathcal H$. For a positive
semidefinite matrix, not necessarily normalized, write
\[
 \ent(P)=-\Tr(P\ln P),\qquad 0\ln0=0.
\]
Thus $\ent$ is entropy in nats on states. Write
\[
 H(q_1,\ldots,q_m)=-\sum_iq_i\ln q_i,\qquad
 h(q)=-q\ln q-(1-q)\ln(1-q)
\]
for the entropy of a probability vector and the binary entropy,
respectively. The determinant entropy function is
\begin{equation}\label{eq:g}
 g(u)=h\!\left(\frac{1+\sqrt{1-4u}}2\right),
 \qquad 0\le u\le\frac14.
\end{equation}
A two-by-two state of determinant $u$ has entropy $g(u)$, by its
characteristic polynomial.

For $\gamma,\nu\in[0,1]$, define the generalized amplitude-damping
channel by
\begin{equation}\label{eq:channel}
 \A_{\gamma,\nu}
 \begin{pmatrix}1-q&z\\\overline z&q\end{pmatrix}
 =\begin{pmatrix}1-t&\sqrt{1-\gamma}\,z\\
 \sqrt{1-\gamma}\,\overline z&t\end{pmatrix},
 \qquad t=(1-\gamma)q+\gamma\nu.
\end{equation}
The input is a state when $0\le q\le1$ and $|z|^2\le q(1-q)$.
Here $\nu$ is the stationary excited-state population; this explicit
convention fixes any possible ambiguity in thermal parameters.

We measure Holevo information and capacity in \emph{bits}:
\begin{equation}\label{eq:holevo}
 \chi(\mathcal N)=\frac1{\ln2}
 \sup_{\{q_a,\rho_a\}}
 \left\{\ent\!\left(\sum_aq_a\mathcal N(\rho_a)\right)
       -\sum_aq_a\ent\bigl(\mathcal N(\rho_a)\bigr)\right\}.
\end{equation}
The supremum runs over all finite ensembles of input states, with
$q_a\ge0$ and $\sum_aq_a=1$. For a tensor-power channel, the states
are unrestricted on the entire input space. An $n$-use classical code assigns a state on the full input space to
each of its $M_n$ equiprobable messages and uses one measurement on all $n$ outputs.
The unassisted classical capacity $C(\mathcal N)$ is the supremum of
rates $R$ for which such codes have average decoding error tending to
zero and $\liminf_{n\to\infty}n^{-1}\log_2 M_n\ge R$.
There is no preshared entanglement or feedback. The channel is
memoryless, so its $n$-use action is $\mathcal N^{\otimes n}$.
The coding theorem gives
\begin{equation}\label{eq:regularization}
 C(\mathcal N)=\sup_{n\ge1}\frac1n\chi(\mathcal N^{\otimes n}).
\end{equation}
Indeed, the mixed-signal coding theorem
\cite{Holevo1998,SchumacherWestmoreland1997} applies to any fixed finite
ensemble of block inputs, treating each block as one letter. It achieves
rates below the ensemble's Holevo information per physical channel use;
see \cite[preprint, Section~2, Eq.~(5)]{Holevo1998} for the finite-alphabet
statement. Conversely, the Holevo bound and Fano's inequality applied to
the uniform output ensemble of any $n$-use code with average error
$\epsilon_n$ give
\[
 (1-\epsilon_n)\log_2 M_n
 \le \chi(\mathcal N^{\otimes n})+\frac{h(\epsilon_n)}{\ln2}.
\]
The normalized asymptotic limit equals the supremum in
\eqref{eq:regularization}: if $a_n=\chi(\mathcal N^{\otimes n})$,
product ensembles give $a_{m+n}\ge a_m+a_n$, and
$0\le a_n\le n\log_2 d_{\mathrm{out}}$, where $d_{\mathrm{out}}$ is
the output dimension. For fixed $k$, writing $n=qk+r$ gives
$a_n\ge q a_k$; taking a lower limit and then the supremum over $k$
proves the assertion. This coding result supplies the operational
interpretation of the entropy optimization below.

For a finite-dimensional channel, define its minimum output entropy in
\emph{nats} by
\begin{equation}\label{eq:min-output}
 S_{\min}(\mathcal N)=\min_{\rho}\ent\bigl(\mathcal N(\rho)\bigr),
\end{equation}
where $\rho$ ranges over all input density matrices. The minimum exists
because the input state space is compact and entropy is continuous in
finite dimensions.

\begin{theorem}\label{thm:main}
For all $\gamma,\nu\in[0,1]$ and every integer $n\ge1$,
\begin{equation}\label{eq:main}
 \chi(\A_{\gamma,\nu}^{\otimes n})
 =n\chi(\A_{\gamma,\nu})
 =\frac n{\ln2}\max_{0\le p\le1}
 \{h((1-\gamma)p+\gamma\nu)-g(v_{\gamma,\nu}(p))\},
\end{equation}
where
\begin{equation}\label{eq:v-main}
 v_{\gamma,\nu}(p)
 =\gamma\nu(1-\nu)+\gamma(1-\gamma)(p-\nu)^2.
\end{equation}
In particular, $C(\A_{\gamma,\nu})=\chi(\A_{\gamma,\nu})$.
If $p$ maximizes \eqref{eq:main}, independent products of the two
equiprobable signals
\begin{equation}\label{eq:pair}
 \phi_\pm=\sqrt{1-p}\,\ket0\ \pm\sqrt p\,\ket1
\end{equation}
attain the Holevo information at every block length.
\end{theorem}

The tensor-power upper bound includes all entangled block signals; a
product ensemble attains it. Equation~\eqref{eq:regularization} turns
this finite-block optimum into an asymptotic reliable rate with
collective decoding.

The rectangular thermal inequality developed below gives the following
stronger consequence.

\begin{corollary}[Additivity with a finite-dimensional partner]
\label{cor:partner-additivity}
Let $\gamma,\nu\in[0,1]$, put $\A=\A_{\gamma,\nu}$, and let
$\mathcal N:\mathcal B(\mathcal H_{\mathrm{in}})\to
\mathcal B(\mathcal H_{\mathrm{out}})$ be completely positive and
trace preserving, where $\mathcal H_{\mathrm{in}}$ and
$\mathcal H_{\mathrm{out}}$ are nonzero finite-dimensional complex
Hilbert spaces. Then
\begin{align}
 \chi(\A\otimes\mathcal N)
 &=\chi(\A)+\chi(\mathcal N),\label{eq:partner-holevo}\\
 S_{\min}(\A\otimes\mathcal N)
 &=S_{\min}(\A)+S_{\min}(\mathcal N),\label{eq:partner-min}\\
 C(\A\otimes\mathcal N)
 &=C(\A)+C(\mathcal N).\label{eq:partner-capacity}
\end{align}
Here $\chi$ is the unconstrained one-shot Holevo capacity in
\eqref{eq:holevo}, and $C$ is the regularized unassisted classical
capacity in \eqref{eq:regularization}.
\end{corollary}

The partner retains its full regularization: $C(\mathcal N)$ can differ
from $\chi(\mathcal N)$.

\subsection{Proof strategy}

For a pure one-qubit input with excitation probability $p$, the output
determinant is $v_{\gamma,\nu}(p)$ and its entropy is
$e(p):=g(v_{\gamma,\nu}(p))$. The central estimate extends this exact
one-site value to every pure $n$-qubit input $\psi$:
\begin{equation}\label{eq:roadmap}
 \ent\bigl(\A_{\gamma,\nu}^{\otimes n}(\proj\psi)\bigr)
 \ge\sum_{j=1}^n e(p_j),
\end{equation}
where $p_j$ is the probability of outcome $1$ at site $j$ in the
computational basis. A bound only on the unconstrained minimum output
entropy would lose these populations and would not suffice for the
Holevo optimization.

To prove \eqref{eq:roadmap}, split the input according to its first
qubit:
$\psi=\sqrt{1-p_1}\ket0\otimes\psi_0+
\sqrt{p_1}\ket1\otimes\psi_1$.
The normalized conditional vectors $\psi_0,\psi_1$ can be nonorthogonal
and can remain entangled on the other sites. For the channel on those
sites choose a Kraus representation $\mathcal N(P)=\sum_J L_JPL_J^*$.
Define column $J$ of $X$ as $L_J\psi_0$ and column $J$ of $Y$ as
$L_J\psi_1$. Their Gram products are the two branch
outputs, and $XY^*$ is the image of the cross operator between the
branches. The thermal block inequality gives
\[
 \ent(\text{full output})\ge e(p_1)
 +(1-p_1)\ent(XX^*)+p_1\ent(YY^*).
\]
This one-step application does not require $\mathcal N$ to be a damping
tensor power: its Kraus columns may have any finite rectangular shape.
For a damping tensor power, convexity of $e$ combines the conditional
populations into the original site populations, and iteration gives
\eqref{eq:roadmap} without restricting the input vector.

The matrix proof leading to this recursion is developed in
\Cref{sec:block,sec:thermal}. For a zero-block factor
$M=\bigl(\begin{smallmatrix}A&B\\ C&0\end{smallmatrix}\bigr)$, with
squared Hilbert--Schmidt masses
$x=\Tr(AA^*)$, $y=\Tr(BB^*)$, $z=\Tr(CC^*)$ summing to one, the entropy
contribution beyond the weighted block entropies is at least $g(yz)$.
We express the entropy deficit as an integral of log-determinant
differences. Joint convexity permits simultaneous averaging over
coordinate signs; this reduces the estimate to a concave scalar
deficit. Zero padding preserves the result for rectangular blocks,
which are required by the Kraus construction. The thermal extension
uses two such factorizations whose products of corner masses agree,
so the same $g$ term survives their convex mixture.

Finally, for an ensemble with mean populations $\overline p_j$, the
entropy of the ensemble-average output is at most
$\sum_j h((1-\gamma)\overline p_j+\gamma\nu)$, by the diagonal-entropy bound in
the computational basis and classical subadditivity. Entropy concavity
and convexity of $e$ first extend \eqref{eq:roadmap} to mixed inputs.
Averaging this bound and using convexity of $e$ bounds the average of
the individual output entropies from below by
$\sum_j e(\overline p_j)$. The difference gives the upper bound in
\Cref{thm:main}; products of the phase pair \eqref{eq:pair} attain it.
For an arbitrary partner, the branch outputs form an ensemble whose
mean is the second output marginal. The entropy of the mean full output
is at most the binary entropy of its first output population plus the
entropy of that same marginal. Subtraction leaves one scalar damping
objective plus at most $(\ln2)\chi(\mathcal N)$. The minimum-output
equality follows from the same one-step estimate, and the capacity
equality follows by iterating against $\mathcal N^{\otimes n}$ before
regularizing.
\Cref{sec:output} proves the recursion and population convexity,
\Cref{sec:holevo} carries out the ensemble optimization, and
\Cref{sec:partners} derives the arbitrary-partner consequences.
\Cref{sec:special} records the ordinary, unital, and endpoint formulas.
\Cref{sec:alternative} gives direct Hessian proofs of the scalar
concavity and log-determinant convexity, providing an independent
analytic route to the same zero-block estimate.

\section{An entropy inequality with one zero block}
\label{sec:block}

Our first objective is to separate the entropy produced by the block
pattern from the entropies within its three nonzero entries. The
rectangular formulation is needed because a channel output and its
Kraus index space can have different dimensions. The following is the
case of the triangular block entropy inequality of Tang, Zhu, Bai,
and Wang \cite[Theorem~4.2]{TangEtAl2026} in which all blocks have the
same rectangular shape, after permuting block rows. We give a direct
proof in the form needed for the thermal extension.

\begin{theorem}\label{thm:block}
Let $A,B,C\in\C^{d\times e}$, and put
\[
 M=\begin{pmatrix}A&B\\ C&0\end{pmatrix},\qquad
 x=\Tr(AA^*),\quad y=\Tr(BB^*),\quad z=\Tr(CC^*).
\]
If $x+y+z=1$, then
\begin{equation}\label{eq:block}
 \ent(MM^*)\ge
 x\ent(AA^*/x)+y\ent(BB^*/y)+z\ent(CC^*/z)+g(yz),
\end{equation}
where each term with zero weight is omitted.
\end{theorem}

The scalar matrix
$\bigl(\begin{smallmatrix}\sqrt x&\sqrt y\\\sqrt z&0\end{smallmatrix}\bigr)$
has squared singular values with sum $1$ and product $yz$.
The term $g(yz)$ is their entropy. To separate this contribution from
the entropies inside the three blocks, write the entropy deficit as
\[
 \Delta=\ent(AA^*)+\ent(BB^*)+\ent(CC^*)-\ent(MM^*).
\]
Since $\ent(AA^*)=x\ent(AA^*/x)-x\ln x$ for $x>0$, and similarly
for the other blocks, the theorem is equivalent to
\[
 \Delta\le H(x,y,z)-g(yz).
\]
We will express $\Delta$ as an integral of log-determinant differences.
Joint convexity then permits averaging over coordinate signs, and
homogeneous scalar concavity combines the resulting coordinates.
This integral and averaging route is also used in
\cite[Appendix~A.2--A.4]{TangEtAl2026}.

Block norm compression provides a related perspective. King
\cite[Theorem~1]{King2003} compared the Schatten norm of a positive
two-by-two block matrix with the norm of the scalar matrix of block
norms. Audenaert \cite[Conjecture~1]{Audenaert2008} studied the
corresponding comparison for arbitrary two-row block matrices and
proved several special cases. The entropy argument below allows
noncommuting blocks and does not require the general norm-compression
conjecture.

\subsection{From entropy to log determinants}

The direct sum of the three individual Gram matrices and the full Gram
matrix have equal trace. The following identity converts their entropy difference
into an integral, including when their dimensions or ranks differ.

\begin{lemma}\label{lem:integral}
If $U,V\ge0$ have equal trace, with possibly different sizes, then
\begin{equation}\label{eq:block-entropy-integral}
 \ent(U)-\ent(V)
 =\int_0^\infty
  \bigl[\ln\det(I+rU)-\ln\det(I+rV)\bigr]\frac{\dd r}{r^2}.
\end{equation}
The integral converges absolutely.
\end{lemma}

\begin{proof}
Let $f(r)$ denote the bracket. Equal trace gives $f(r)=O(r^2)$
at zero, while $f(r)=O(\ln r)$ at infinity. Hence the integral
converges absolutely and $f(r)/r$ vanishes at both endpoints.
Integration by parts reduces it to $\int_0^\infty f'(r)\dd r/r$.
If $a_i$ and $b_j$ are the positive eigenvalues of $U$ and $V$, then
\[
 \frac{f'(r)}r
 =\sum_i\frac{a_i}{r(1+ra_i)}
  -\sum_j\frac{b_j}{r(1+rb_j)}.
\]
A primitive of the term for $a>0$ is
$a\ln\bigl(r/(1+ra)\bigr)$. In the signed sum, its lower endpoint
limit is zero because the coefficients of $\ln r$ cancel by equal
trace. Its upper endpoint contribution is $-a\ln a$.
Taking their difference gives $\ent(U)-\ent(V)$. Zero eigenvalues
contribute nothing.
\end{proof}

\subsection{Joint convexity by a contraction}

The log-determinant function used in the averaging step is the following.
Its joint convexity follows by specialization and continuity from Hiai's
trace-function theorem \cite[preprint, Theorem~5.2]{Hiai2016}.
We give a direct proof: the contraction argument is the specialization
of \cite[Lemma~A.1]{TangEtAl2026} to a two-by-two block matrix.

\begin{lemma}\label{lem:logdet}
For each fixed $Z\in\C^{N\times N}$, the function
\begin{equation}\label{eq:block-G}
 G_Z(D,E)=\ln\det(D+ZE^{-1}Z^*)-\ln\det D,
 \qquad D,E>0,
\end{equation}
is jointly convex on pairs of positive definite Hermitian matrices.
\end{lemma}

\begin{proof}
Set
\[
 \mathcal D=\begin{pmatrix}D&0\\0&E\end{pmatrix},\qquad
 \mathcal C=\begin{pmatrix}D&Z\\-Z^*&E\end{pmatrix}.
\]
The block determinant identity gives
\[
 \det\mathcal C=\det E\det(D+ZE^{-1}Z^*)>0,
 \qquad G_Z(D,E)=\ln\det\mathcal C-\ln\det\mathcal D.
\]
Although $\mathcal C$ is not generally Hermitian, its determinant is
positive and its logarithm here is the real logarithm.
Along an affine line in $(D,E)$, let $\mathcal H$ be the derivative
of $\mathcal D$, which is also the derivative of $\mathcal C$.
The determinant derivative formula gives
\begin{equation}\label{eq:alt-second}
 G_Z''=\Tr(\mathcal D^{-1}\mathcal H\mathcal D^{-1}\mathcal H)
       -\Tr(\mathcal C^{-1}\mathcal H\mathcal C^{-1}\mathcal H).
\end{equation}
At the point under consideration put
\[
 L=\mathcal D^{-1/2}\mathcal H\mathcal D^{-1/2},\qquad
 K=\mathcal D^{-1/2}(\mathcal C-\mathcal D)\mathcal D^{-1/2}.
\]
Then $L=L^*$ and $K^*=-K$. Consequently
$\|(I+K)v\|_2^2=\|v\|_2^2+\|Kv\|_2^2$ for every vector $v$,
and $R=(I+K)^{-1}$ exists with operator norm at most one.
Cyclicity of trace rewrites \eqref{eq:alt-second} as
\[
 G_Z''=\|L\|_{\mathrm{HS}}^2-\Tr((RL)^2).
\]
The last trace is real because both log determinants along the line
are real. For any complex matrix $B$, entrywise Cauchy--Schwarz gives
$|\Tr(B^2)|=|\sum_{i,j}B_{ij}B_{ji}|\le\|B\|_{\mathrm{HS}}^2$.
Therefore
\[
 \Tr((RL)^2)\le|\Tr((RL)^2)|
 \le\|RL\|_{\mathrm{HS}}^2\le\|L\|_{\mathrm{HS}}^2.
\]
Every second directional derivative is thus nonnegative. The positive
definite domain is convex, so $G_Z$ is convex along each line segment,
as required.
\end{proof}

\subsection{The scalar deficit and its homogeneity}

We next use the same convexity to control the scalar entropy deficits
that sign averaging will produce. For $x,y,z\ge0$, define
\begin{align}
 m&=x+y+z,\qquad
 \lambda_\pm=\frac{m\pm\sqrt{m^2-4yz}}2,
 \label{eq:block-lambda}\\
 F(x,y,z)&=-x\ln x-y\ln y-z\ln z
           +\lambda_+\ln\lambda_++\lambda_-\ln\lambda_-.
 \label{eq:block-F}
\end{align}
Here $m^2\ge4yz$, and $\lambda_\pm$ are nonnegative with sum $m$
and product $yz$. Thus $F$ is continuous on the nonnegative octant,
with the usual convention at zero. On the probability simplex it is
exactly the deficit in our target: $F(x,y,z)=H(x,y,z)-g(yz)$.

\begin{lemma}\label{lem:scalar}
The function $F$ is positively homogeneous of degree one and concave
on the probability simplex. For every finite collection of nonnegative
triples,
\begin{equation}\label{eq:block-superadditive}
 \sum_i F(x_i,y_i,z_i)
 \le F\!\left(\sum_i x_i,\sum_i y_i,\sum_i z_i\right).
\end{equation}
\end{lemma}

\begin{proof}
The nonnegative entries in the two lists $(x,y,z)$ and
$(\lambda_+,\lambda_-)$ have the same sum. Applying
\Cref{lem:integral} to the corresponding diagonal matrices gives
\begin{align}
 F(x,y,z)
 &=\int_0^\infty
 \ln\frac{(1+rx)(1+ry)(1+rz)}
          {(1+r\lambda_+)(1+r\lambda_-)}\frac{\dd r}{r^2}\notag\\
 &=\int_0^\infty\left[\ln(1+rx)
 -\ln\!\left(1+\frac{rx}{(1+ry)(1+rz)}\right)\right]
          \frac{\dd r}{r^2}.
 \label{eq:scalar-deficit-integral}
\end{align}
The second equality uses
$(1+r\lambda_+)(1+r\lambda_-)=1+r(x+y+z)+r^2yz$;
the integral is absolutely convergent by the lemma.
For $x=1$ its bracket is
\[
 \ln(1+r)-G_{\sqrt r}(1+ry,1+rz).
\]
By \Cref{lem:logdet}, this is concave in $(y,z)$ for every $r>0$.
Integrating the concavity inequality shows that $F(1,\cdot,\cdot)$
is concave on the nonnegative quadrant.

Under scaling of $(x,y,z)$ by a positive constant, both
$\lambda_\pm$ scale by the same constant. Their sum is $m$, so the
logarithms of the scaling constant cancel in \eqref{eq:block-F}.
This proves homogeneity, including scale zero by continuity.
For a nonempty collection with all $x_i>0$, set $x=\sum_i x_i$.
Homogeneity and the preceding concavity now give
\begin{align*}
 \sum_iF(x_i,y_i,z_i)
 &=x\sum_i\frac{x_i}{x}F(1,y_i/x_i,z_i/x_i)\\
 &\le xF\!\left(1,\frac{\sum_i y_i}{x},
                     \frac{\sum_i z_i}{x}\right)
 =F\!\left(x,\sum_i y_i,\sum_i z_i\right).
\end{align*}
For a collection containing zero first coordinates, replace every
$x_i$ by $x_i+\varepsilon$, apply this inequality, and let
$\varepsilon\downarrow0$. Continuity of $F$ gives
\eqref{eq:block-superadditive}, even when every $x_i=0$.
The empty collection gives zero on both sides. Finally, applying
superadditivity to the two triples $\theta a$ and $(1-\theta)b$
and using homogeneity gives
$F(\theta a+(1-\theta)b)\ge\theta F(a)+(1-\theta)F(b)$
for nonnegative triples $a,b$ and $0\le\theta\le1$.
In particular, $F$ is concave on the probability simplex.
\end{proof}

\subsection{Sign averaging and the block estimate}

The two convexity consequences are now in place. The Schur complement
will put the matrix deficit into the form involving $G_Z$; averaging
over signs will replace its matrix arguments by diagonal entries;
and \Cref{lem:scalar} will combine the resulting scalar deficits.

\begin{proof}[Proof of \Cref{thm:block}]
First suppose that all three blocks are square of size $N$, and
temporarily allow arbitrary nonnegative masses $x,y,z$.
The matrices
\[
 U=AA^*\oplus BB^*\oplus CC^*,\qquad V=MM^*
\]
have equal trace. For $r>0$, put $D=I+rBB^*$ and $E=I+rC^*C$.
The Schur complement of $I+rCC^*$ in $I+rMM^*$ yields
\begin{equation}\label{eq:block-schur}
 \det(I+rMM^*)
 =\det(I+rCC^*)\det(D+rAE^{-1}A^*),
\end{equation}
where
$I-rC^*(I+rCC^*)^{-1}C=(I+rC^*C)^{-1}$.
Consequently the integrand bracket in \eqref{eq:block-entropy-integral}
is
\begin{equation}\label{eq:block-integrand}
 \ln\det(I+rAA^*)-G_{\sqrt r A}(D,E).
\end{equation}

Choose a singular-value decomposition $A=U_A\Sigma V_A^*$.
Replacing $(A,B,C)$ by $(\Sigma,U_A^*B,CV_A)$ amounts to
multiplying $M$ on the left and right by block-diagonal unitaries;
all entropies and masses are preserved. In these coordinates write
\[
 A=\diag(\sqrt{x_1},\ldots,\sqrt{x_N}),\qquad
 y_i=(BB^*)_{ii},\qquad z_i=(C^*C)_{ii}.
\]
These entries are nonnegative and satisfy
$\sum_i(x_i,y_i,z_i)=(x,y,z)$.
For each $\varepsilon\in\{-1,1\}^N$, let
$S_\varepsilon=\diag(\varepsilon_1,\ldots,\varepsilon_N)$.
Because $S_\varepsilon A S_\varepsilon=A$,
\[
 G_{\sqrt r A}(S_\varepsilon DS_\varepsilon,
                S_\varepsilon ES_\varepsilon)
 =G_{\sqrt r A}(D,E).
\]
Averaging simultaneously over all signs removes the off-diagonal
entries of both arguments. Joint convexity from \Cref{lem:logdet}
therefore gives
\[
 G_{\sqrt r A}(D,E)\ge
 G_{\sqrt r A}\bigl(\diag(1+ry_i),\diag(1+rz_i)\bigr).
\]
Because $G$ is subtracted in \eqref{eq:block-integrand}, that bracket
is bounded above by
\begin{align}
 &\sum_i\left[\ln(1+rx_i)
 -\ln\!\left(1+\frac{rx_i}{(1+ry_i)(1+rz_i)}\right)\right]\notag\\
 &\hspace{12mm}=
 \sum_i\ln\frac{(1+rx_i)(1+ry_i)(1+rz_i)}
                    {1+r(x_i+y_i+z_i)+r^2y_iz_i}.
 \label{eq:block-scalar-integrands}
\end{align}
By \eqref{eq:scalar-deficit-integral}, each summand integrates against
$\dd r/r^2$ to $F(x_i,y_i,z_i)$, and these finitely many integrals
converge absolutely. Integrating
the inequality and applying \Cref{lem:scalar} proves
\begin{equation}\label{eq:block-unnormalized}
 \ent(AA^*)+\ent(BB^*)+\ent(CC^*)-\ent(MM^*)
 \le\sum_iF(x_i,y_i,z_i)\le F(x,y,z).
\end{equation}
When $x+y+z=1$, the last expression is $H(x,y,z)-g(yz)$.
For $x>0$,
$\ent(AA^*)=x\ent(AA^*/x)-x\ln x$, and the same identity holds
for the other two blocks. Substitution gives \eqref{eq:block}.
If a mass vanishes, the corresponding block is zero, and its
contribution is zero. Thus the proof includes every boundary case;
all inverses used above are of positive definite matrices even
when $A,B,C$ are singular.

Finally, for $d\times e$ blocks let $N=\max(d,e)$ and choose the
coordinate isometries $J:\C^d\to\C^N$ and $K:\C^e\to\C^N$.
Replace each block $A$ by $JAK^*$, and likewise for $B,C$.
The resulting square block matrix is
$(J\oplus J)M(K\oplus K)^*$, so this operation only adds zero
singular values to $M$ and to each block. All masses and entropies
in \eqref{eq:block} are unchanged. The square case therefore proves
the rectangular case as well.
\end{proof}

\subsection{A scalar convexity fact for the channel application}

The scalar deficit estimate completes the zero-block theorem. The later
channel argument also needs a different property of the same entropy
function: its convexity as a function of the square root of the
determinant. We record it separately.

\begin{lemma}\label{lem:binary-spectrum}
The function $s\mapsto g(s^2)$ is nondecreasing and convex on $[0,1/2]$.
Consequently, $g$ is nondecreasing on $[0,1/4]$.
\end{lemma}

\begin{proof}
For $0<u<1/4$, put
\[
 t=\sqrt{1-4u},\qquad L=\frac{\atanh t}{t},\qquad
 \atanh t=\frac12\ln\frac{1+t}{1-t}.
\]
Differentiating the definition of $g$ gives
\begin{equation}\label{eq:g-derivatives}
 g'(u)=2L,\qquad
 g''(u)=\frac{4L-1/u}{t^2},\qquad
 ug''(u)+\frac12g'(u)=\frac{L-1}{t^2}\ge0.
\end{equation}
The last inequality follows from $\atanh t\ge t$, since
$(\atanh t)'=(1-t^2)^{-1}\ge1$ and both functions vanish at zero.
For $0<s<1/2$, these identities give
\[
 \frac{\mathrm d}{\mathrm ds}g(s^2)=2s g'(s^2)\ge0,\qquad
 \frac{\mathrm d^2}{\mathrm ds^2}g(s^2)
 =4\left(s^2g''(s^2)+\frac12g'(s^2)\right)\ge0.
\]
Continuity supplies both conclusions at the endpoints; no boundary
derivatives are required. Since $u\mapsto\sqrt u$ is nondecreasing,
$g(u)=g((\sqrt u)^2)$ is nondecreasing as well.
\end{proof}

The function $\kappa\mapsto g(\kappa^2/4)/\ln2$ is the
entropy--concurrence function of Wootters
\cite[Eq.~(8)]{Wootters1998}. Its monotonicity and convexity are the
scalar properties just proved. We use these properties, rather than
any additivity statement about entanglement of formation.

\section{A thermal block entropy inequality}\label{sec:thermal}

The zero-block theorem does not apply directly when thermal excitation
and decay both occur. We will express the resulting state as a convex
mixture of two states to which it does apply, preserving the same scalar
entropy term in each. We first recall two elementary entropy facts,
including their proofs to fix the directions of the inequalities used
below. For a state $P$,
\begin{equation}\label{eq:diagonal}
 \ent(P)\le H(\diag P).
\end{equation}
Indeed, its diagonal is obtained from its eigenvalue vector by the
doubly stochastic matrix of squared absolute unitary entries.
Concavity of $-s\ln s$, applied row by row and then summed, gives
\eqref{eq:diagonal}. Entropy is also concave on states:
\begin{equation}\label{eq:concavity}
 \ent\!\left(\sum_iw_iP_i\right)\ge\sum_iw_i\ent(P_i).
\end{equation}
To see this, choose a basis diagonalizing the mixture, apply classical
entropy concavity in that basis, and then use \eqref{eq:diagonal} on
each $P_i$.

The next result is the extension needed for thermal channels.

\begin{theorem}\label{thm:thermal}
Let $X,Y$ be complex matrices of the same rectangular shape with
$\Tr(XX^*)=\Tr(YY^*)=1$. Suppose
$\alpha,\beta,\eta,\delta\ge0$ sum to one and
\begin{equation}\label{eq:thermal-condition}
 \alpha\delta-\beta\eta=K^2,\qquad K\ge0.
\end{equation}
Set
\begin{equation}\label{eq:thermal-state}
 T=\begin{pmatrix}
 \alpha XX^*+\beta YY^*&KXY^*\\
 KYX^*&\eta XX^*+\delta YY^*
 \end{pmatrix},\qquad
 U_0=\alpha+\beta,\quad U_1=\eta+\delta,
\end{equation}
and $u=U_0U_1-K^2$. Then $T$ is a state, $0\le u\le1/4$, and
\begin{equation}\label{eq:thermal-bound}
 \ent(T)\ge(\alpha+\eta)\ent(XX^*)
             +(\beta+\delta)\ent(YY^*)+g(u).
\end{equation}
\end{theorem}

\begin{proof}
The coefficient identity gives
\begin{equation}\label{eq:thermal-u}
 u=\beta U_1+\eta U_0.
\end{equation}
It follows that $0\le u\le U_0U_1\le1/4$.

Suppose first that $u>0$, so $U_0,U_1>0$. To construct the mixture,
keep $X,Y,K$ and the row sums $U_0,U_1$ fixed. Writing
$\alpha=U_0-\beta$ and $\delta=U_1-\eta$, the coefficient condition
\eqref{eq:thermal-condition} becomes the line equation
$U_1\beta+U_0\eta=u$. Its nonnegative solutions form the segment
joining $(\beta,\eta)=(u/U_1,0)$ and $(0,u/U_0)$. At either endpoint
one coefficient vanishes, permitting a zero-block Gram factorization.
The fixed row sums and cross coefficient retain the same value of $u$.
These endpoints and the weights expressing the original coefficients
as their convex mixture are
\begin{equation}\label{eq:thermal-table}
 \begin{array}{c|cccc|c}
 i&\alpha_i&\beta_i&\eta_i&\delta_i&w_i\\ \hline
 1&K^2/U_1&u/U_1&0&U_1&\beta U_1/u\\[2pt]
 2&U_0&0&u/U_0&K^2/U_0&\eta U_0/u
 \end{array}
\end{equation}
Both coefficient sets are nonnegative and have the original row sums
$U_0,U_1$, since $K^2+u=U_0U_1$. Each therefore sums to one and
satisfies $\alpha_i\delta_i=K^2$. Moreover, $w_1+w_2=1$ by
\eqref{eq:thermal-u}, and
\[
 \sum_iw_i\beta_i=\beta,\qquad
 \sum_iw_i\eta_i=\eta.
\]
The fixed row sums then recover $\alpha$ and $\delta$ as well.

Let $T_i$ denote \eqref{eq:thermal-state} with the $i$th coefficient
set. Both $T_i$ have cross block $KXY^*$, so
$T=w_1T_1+w_2T_2$. They have Gram factorizations
\begin{equation}\label{eq:thermal-factors}
 \begin{split}
 M_1&=\begin{pmatrix}
 \sqrt{\alpha_1}X&\sqrt{\beta_1}Y\\
 \sqrt{\delta_1}Y&0
 \end{pmatrix},\\
 M_2&=\begin{pmatrix}
 \sqrt{\alpha_2}X&0\\
 \sqrt{\delta_2}Y&\sqrt{\eta_2}X
 \end{pmatrix},\qquad T_i=M_iM_i^*.
 \end{split}
\end{equation}
In particular, $T_i$ and $T$ are positive semidefinite of trace one.
Apply \Cref{thm:block} to $M_1$ and to $M_2$ after swapping its two
block rows. The swap is unitary on the output space and preserves
entropy. The respective products of the off-corner block masses are
\[
 \beta_1\delta_1=u,\qquad \eta_2\alpha_2=u.
\]
Thus \emph{the same} scalar term occurs in both inequalities:
\[
 \ent(T_i)\ge(\alpha_i+\eta_i)\ent(XX^*)
                 +(\beta_i+\delta_i)\ent(YY^*)+g(u).
\]
Concavity \eqref{eq:concavity} and recovery of the original
coefficients now give \eqref{eq:thermal-bound}. Zero weights in
\eqref{eq:thermal-table} do not change this argument.

It remains to consider $u=0$. If a row sum vanishes, all its
nonnegative coefficients vanish and $K=0$. The other diagonal block
is a convex mixture of $XX^*$ and $YY^*$, so positivity, normalization,
and \eqref{eq:thermal-bound} follow from \eqref{eq:concavity} and
$g(0)=0$. If both row sums are positive, \eqref{eq:thermal-u} forces
$\beta=\eta=0$, and $K=\sqrt{\alpha\delta}$. Then the first
factorization in \eqref{eq:thermal-factors}, with coefficients
$(\alpha,0,0,\delta)$, applies directly. It proves positivity and
reduces \eqref{eq:thermal-bound} to \Cref{thm:block} with a zero block
mass. These cases exhaust the possibilities.
\end{proof}

\begin{remark}\label{rem:thermal-sharp}
No orthogonality condition on $X,Y$ and no commutativity of their Gram
products is used. The bound is sharp when $X=Y$: in that case
$T=\bigl(\begin{smallmatrix}U_0&K\\K&U_1\end{smallmatrix}\bigr)
\otimes XX^*$, whose entropy is $g(u)+\ent(XX^*)$.
\end{remark}

\section{From the thermal inequality to channel outputs}\label{sec:output}

We first turn the thermal matrix inequality into an entropy bound for
$\A_{\gamma,\nu}\otimes\mathcal N$, with any finite-dimensional
partner channel $\mathcal N$. Its two matrices will encode the outputs
of conditional input branches; a common Kraus index will preserve their
cross operator. We then iterate this bound for repeated damping uses,
retaining the input population at every site.

\subsection{The one-qubit entropy and its convexity}

Fix $\gamma,\nu\in[0,1]$ and abbreviate
$\A=\A_{\gamma,\nu}$. Put
\begin{equation}\label{eq:coefficients}
 k=\sqrt{1-\gamma},\qquad b=\gamma(1-\nu),\qquad
 c=\gamma\nu,\qquad a=1-c,\qquad d=1-b.
\end{equation}
Then
\begin{equation}\label{eq:coefficient-identities}
 a+c=b+d=1,\qquad ad-bc=k^2.
\end{equation}
The population transition matrix is
$\bigl(\begin{smallmatrix}a&b\\c&d\end{smallmatrix}\bigr)$ and
the coherence multiplier is $k$. A Kraus representation is
\begin{equation}\label{eq:kraus}
 \begin{aligned}
 L_0&=\sqrt{1-\nu}\,\diag(1,k),&
 L_1&=\sqrt b\,\ket0\bra1,\\
 L_2&=\sqrt\nu\,\diag(k,1),&
 L_3&=\sqrt c\,\ket1\bra0.
 \end{aligned}
\end{equation}
Direct multiplication gives $\A(P)=\sum_{i=0}^3L_iPL_i^*$ and
$\sum_{i=0}^3L_i^*L_i=I$, including all parameter endpoints.

For $p\in[0,1]$, define
\begin{equation}\label{eq:single-entropy}
 \begin{split}
 v(p)&=(a(1-p)+bp)(c(1-p)+dp)-k^2p(1-p)\\
     &=ac(1-p)^2+bdp^2+2bcp(1-p),\qquad e(p)=g(v(p)).
 \end{split}
\end{equation}
Any pure qubit state of excitation probability $p$ has coherence
modulus $\sqrt{p(1-p)}$. Hence $v(p)$ is its output determinant and
$e(p)$ its output entropy. In particular, $v(p)\in[0,1/4]$.
Expanding \eqref{eq:coefficients} in \eqref{eq:single-entropy} gives
\eqref{eq:v-main}. If $\gamma>0$ and $0<\nu<1$,
then $v(p)\ge\gamma\nu(1-\nu)>0$ for every $p$. Every pure input
therefore has a positive definite output. A mixed input is a convex
combination of pure inputs, so its output is positive definite as well.
Thus these thermal parameters lie outside the pure-output class treated
by Tang, Zhu, Bai, and Wang
\cite[Corollary~5.1 and Lemma~5.2]{TangEtAl2026}.

\begin{lemma}\label{lem:e-convex}
For every $\gamma,\nu\in[0,1]$, the function
$p\mapsto g(\gamma\nu(1-\nu)+\gamma(1-\gamma)(p-\nu)^2)$
is convex on $[0,1]$.
\end{lemma}
\begin{proof}
The completed-square formula \eqref{eq:v-main} gives
\[
 \sqrt{v(p)}=
 \left\|\begin{pmatrix}
 \sqrt{\gamma\nu(1-\nu)}\\
 \sqrt{\gamma(1-\gamma)}\,(p-\nu)
 \end{pmatrix}\right\|_2.
\]
This is a convex function of $p$, being the norm of an affine vector,
even when either coefficient vanishes. By
\Cref{lem:binary-spectrum}, $s\mapsto g(s^2)$ is nondecreasing and convex on
$[0,1/2]$, which contains the range of $\sqrt{v}$. A nondecreasing
convex function of a convex function is convex: apply the inner
convexity, then monotonicity, then the outer convexity.
This proves the assertion, with no parameter exclusions.
\end{proof}

\subsection{A single damping use with an arbitrary partner}

The following estimate is the channel form of \Cref{thm:thermal}.
Its scalar term is the pure one-qubit entropy $e(p)$, and the remaining
terms are the output entropies of the two conditional branches.

\begin{lemma}\label{lem:partner-branch}
Let $\gamma,\nu\in[0,1]$, put $\A=\A_{\gamma,\nu}$ and
$e(p)=g(v_{\gamma,\nu}(p))$, and let
$\mathcal N:\mathcal B(\mathcal H_{\mathrm{in}})\to
\mathcal B(\mathcal H_{\mathrm{out}})$ be completely positive and
trace preserving on nonzero finite-dimensional complex Hilbert spaces.
For a unit vector $\psi\in\C^2\otimes\mathcal H_{\mathrm{in}}$, write
\[
 \psi=\sqrt{1-p}\,\ket0\otimes\psi_0
       +\sqrt p\,\ket1\otimes\psi_1,
 \qquad 0\le p\le1,
\]
where each $\psi_i$ is a unit vector and a zero-weight branch may be
chosen arbitrarily. Then
\begin{equation}\label{eq:partner-branch}
 \begin{split}
 \ent\bigl((\A\otimes\mathcal N)(\proj\psi)\bigr)
 &\ge e(p)+(1-p)\ent\bigl(\mathcal N(\proj{\psi_0})\bigr)\\
 &\quad+p\ent\bigl(\mathcal N(\proj{\psi_1})\bigr).
 \end{split}
\end{equation}
\end{lemma}

\begin{proof}
Choose a finite Kraus representation
\[
 \mathcal N(P)=\sum_{\ell=1}^r V_\ell P V_\ell^*,
 \qquad \sum_{\ell=1}^r V_\ell^*V_\ell=I.
\]
Here $V_\ell:\mathcal H_{\mathrm{in}}\to\mathcal H_{\mathrm{out}}$.
Complete positivity in finite dimensions gives such a finite family;
the last identity is trace preservation. If
$d_{\mathrm{out}}=\dim\mathcal H_{\mathrm{out}}$, form
$X,Y\in\C^{d_{\mathrm{out}}\times r}$ by giving their $\ell$th columns
the values $V_\ell\psi_0$ and $V_\ell\psi_1$, respectively. The common
column index gives the three Gram identities and normalizations
\begin{equation}\label{eq:partner-grams}
 \begin{split}
 XX^*&=\mathcal N(\proj{\psi_0}),\qquad
 YY^*=\mathcal N(\proj{\psi_1}),\\
 XY^*&=\mathcal N(\ket{\psi_0}\bra{\psi_1}),\qquad
 \Tr(XX^*)=\Tr(YY^*)=1.
 \end{split}
\end{equation}
No equality between $d_{\mathrm{out}}$, $r$, and the input dimension
is required, and the branch vectors need not be orthogonal.

Set $x=1-p$ and $y=p$, and use $a,b,c,d,k$ from
\eqref{eq:coefficients}. Applying the two channels to the branch
decomposition gives
\[
 (\A\otimes\mathcal N)(\proj\psi)=
 \begin{pmatrix}
 axXX^*+byYY^*&k\sqrt{xy}\,XY^*\\
 k\sqrt{xy}\,YX^*&cxXX^*+dyYY^*
 \end{pmatrix}.
\]
This is \eqref{eq:thermal-state} with
\begin{equation}\label{eq:channel-substitution}
 \alpha=ax,\quad\beta=by,\quad\eta=cx,\quad\delta=dy,
 \qquad K=k\sqrt{xy}.
\end{equation}
The coefficients are nonnegative, sum to one, and satisfy
$\alpha\delta-\beta\eta=(ad-bc)xy=K^2$.
Their column sums are $\alpha+\eta=x$ and $\beta+\delta=y$.
Moreover, their scalar parameter is
\[
 u=(ax+by)(cx+dy)-k^2xy=v(p).
\]
Thus \Cref{thm:thermal} applies to the arbitrary rectangular matrices
$X,Y$ and gives \eqref{eq:partner-branch}. The theorem includes $u=0$,
and the zero-weight branch terms vanish, so this argument includes
$p=0,1$ and every boundary value of $\gamma,\nu$.
\end{proof}

\subsection{Iteration over a block}

We now take the partner to be the channel on the remaining qubits.
Convexity of $e$ will combine the conditional populations at each later
site into that site's original population.

\begin{proposition}\label{prop:pure-bound}
For every $\gamma,\nu\in[0,1]$, put $\A=\A_{\gamma,\nu}$ and
$e(p)=g(v_{\gamma,\nu}(p))$. For every integer $n\ge1$ and
every unit vector $\psi$ on $n$ qubits,
let $p_j$ be the probability of outcome $1$ at site $j$ in the
computational basis. Then
\begin{equation}\label{eq:pure-bound}
 \ent\bigl(\A^{\otimes n}(\proj\psi)\bigr)
 \ge\sum_{j=1}^n e(p_j).
\end{equation}
\end{proposition}

\begin{proof}
For $n=1$ there is equality by \eqref{eq:single-entropy}. Suppose the
claim holds for $n-1$ and write an arbitrary pure input as
\begin{equation}\label{eq:branch-decomposition}
 \psi=\sqrt x\,\ket0\otimes\psi_0
          +\sqrt y\,\ket1\otimes\psi_1,
 \qquad x=1-p_1,\quad y=p_1.
\end{equation}
The branch vectors are normalized but need not be orthogonal. For a
zero-weight branch choose any unit vector. The decomposition is valid
for every input vector, and each conditional vector may be entangled
across the remaining sites.

Apply \Cref{lem:partner-branch} with
$\mathcal N=\A^{\otimes(n-1)}$, and write
$\sigma_i=\A^{\otimes(n-1)}(\proj{\psi_i})$ for the two branch
outputs. The lemma gives the explicit recursion
\begin{equation}\label{eq:recursion}
 \ent\bigl(\A^{\otimes n}(\proj\psi)\bigr)
 \ge e(p_1)+x\ent(\sigma_0)+y\ent(\sigma_1).
\end{equation}
Apply the induction hypothesis to both pure branch inputs. If
$p_{j|i}$ is the excitation probability of site $j>1$ in $\psi_i$,
then
\[
 p_j=xp_{j|0}+yp_{j|1}.
\]
This identity follows from $\langle0|1\rangle=0$ on the first site;
it does not require $\psi_0\perp\psi_1$. By \Cref{lem:e-convex},
\[
 x\ent(\sigma_0)+y\ent(\sigma_1)
 \ge\sum_{j=2}^n\bigl(xe(p_{j|0})+ye(p_{j|1})\bigr)
 \ge\sum_{j=2}^ne(p_j).
\]
Together with \eqref{eq:recursion}, this proves the induction step.
All terms from zero-weight branches have zero coefficient, so the
argument also includes $p_1=0,1$.
\end{proof}

The same population bound holds for mixed inputs. This form will let us
bound every signal in a Holevo ensemble directly.

\begin{corollary}\label{cor:mixed-bound}
For every $\gamma,\nu\in[0,1]$, integer $n\ge1$, and density matrix
$\rho$ on $(\C^2)^{\otimes n}$, let $p_j$ be its excitation
probability at site $j$. Then
\[
 \ent\bigl(\A_{\gamma,\nu}^{\otimes n}(\rho)\bigr)
 \ge\sum_{j=1}^n g(v_{\gamma,\nu}(p_j)).
\]
\end{corollary}
\begin{proof}
Take a finite pure-state decomposition $\rho=\sum_aq_a\proj{\psi_a}$,
and let $p_{a,j}$ be the corresponding populations. Entropy concavity,
\Cref{prop:pure-bound}, and convexity from \Cref{lem:e-convex} give
\[
 \ent\bigl(\A^{\otimes n}(\rho)\bigr)
 \ge\sum_aq_a\sum_j e(p_{a,j})
 \ge\sum_j e\!\left(\sum_aq_ap_{a,j}\right)=\sum_j e(p_j).
\]
\end{proof}

\section{The ensemble bound and its attainment}\label{sec:holevo}

We now turn the population-dependent entropy bound into the full
Holevo assertion. This step is essential: an unconstrained
minimum-output-entropy bound alone would not give the result.

\begin{proof}[Proof of \Cref{thm:main}]
Fix any finite ensemble $\{q_a,\rho_a\}$ for $n$ uses. Each $\rho_a$
is an arbitrary density matrix on all $n$ input qubits and may be
mixed or entangled across sites. Put
\[
 p_{a,j}=\text{excitation probability at site $j$ in $\rho_a$},
 \qquad \overline p_j=\sum_aq_ap_{a,j}.
\]
The mixed-input bound in \Cref{cor:mixed-bound} and the convexity of
$e$ from \Cref{lem:e-convex} give
\begin{equation}\label{eq:ensemble-lower}
 \sum_aq_a\ent\bigl(\A^{\otimes n}(\rho_a)\bigr)
 \ge\sum_{j=1}^n\sum_aq_ae(p_{a,j})
 \ge\sum_{j=1}^ne(\overline p_j).
\end{equation}

Let $\overline\sigma=\sum_aq_a\A^{\otimes n}(\rho_a)$.
By \eqref{eq:diagonal}, its entropy is at most the Shannon entropy
of its joint computational-basis output distribution. The marginal
probability of outcome $1$ at site $j$ is
$t_j=(1-\gamma)\overline p_j+\gamma\nu$: the other channel factors
are trace preserving, so the reduced output is the one-site channel
applied to the reduced input. For any joint distribution, the chain
rule and concavity give
\[
 H(Z_1,\ldots,Z_n)
 =\sum_{j=1}^nH(Z_j\mid Z_1,\ldots,Z_{j-1})
 \le\sum_{j=1}^nH(Z_j).
\]
Here the chain rule follows by factoring each nonzero joint
probability into conditional probabilities, and each average
conditional entropy is at most its marginal entropy by concavity.
Consequently
\begin{equation}\label{eq:ensemble-upper}
 \ent(\overline\sigma)\le
 \sum_{j=1}^nh((1-\gamma)\overline p_j+\gamma\nu).
\end{equation}
Subtracting \eqref{eq:ensemble-lower} from
\eqref{eq:ensemble-upper} and taking the supremum over ensembles yields
\begin{equation}\label{eq:holevo-upper}
 (\ln2)\chi(\A^{\otimes n})
 \le n\max_{0\le p\le1}
       \{h((1-\gamma)p+\gamma\nu)-e(p)\}.
\end{equation}

For the matching lower bound, the continuous objective attains its
maximum on $[0,1]$. Choose any maximizer $p$ and use the equiprobable
pair \eqref{eq:pair}. Their individual output entropies equal $e(p)$;
their opposite coherences cancel in the
average output, which is diagonal with excited-state population
$(1-\gamma)p+\gamma\nu$. Hence this ensemble attains the right-hand
one-use objective exactly for every $\nu$.

On $n$ uses take the $2^n$ products of these vectors, with equal
probabilities $2^{-n}$. Each individual output and their average are
tensor products. Entropy is additive on such products by their product
eigenvalues, so the resulting finite ensemble attains $n$ times the
one-use objective. At $p=0$ or $1$, repeated labels can simply be
merged. This attains \eqref{eq:holevo-upper} at every $n$. Substitution
of \eqref{eq:single-entropy} and \eqref{eq:v-main} proves
\eqref{eq:main}, and \eqref{eq:regularization} gives the capacity
statement.
\end{proof}

\section{Additivity with a finite-dimensional partner}\label{sec:partners}

The branch estimate in \Cref{lem:partner-branch} applies to every
finite-dimensional partner. We now use it to prove the three assertions
of \Cref{cor:partner-additivity}. For the Holevo bound, the key point is
that the conditional partner outputs form an ensemble whose mean is
exactly the partner marginal of the mean full output.

\begin{proof}[Proof of Corollary~\ref{cor:partner-additivity}]
For the Holevo upper bound, it suffices to consider pure-state ensembles.
Every mixed input has a finite spectral decomposition. Refining each
signal into its pure components leaves the mean output unchanged and,
by entropy concavity \eqref{eq:concavity}, does not increase the mean
individual output entropy. Thus refinement cannot decrease Holevo
information; see also Schumacher--Westmoreland
\cite{SchumacherWestmoreland1997}.
The refined pure inputs may still be entangled between the two factors.

Fix a finite ensemble $\{q_a,\proj{\psi_a}\}$ of such inputs.
For each signal use
the decomposition in Lemma~\ref{lem:partner-branch}, with population $p_a$
and branch vectors $\psi_{a,0},\psi_{a,1}$. Put
\[
 \begin{gathered}
 \overline p=\sum_aq_ap_a,\qquad
 w_{a,0}=q_a(1-p_a),\qquad w_{a,1}=q_ap_a,\\
 \tau_{a,i}=\mathcal N(\proj{\psi_{a,i}}).
 \end{gathered}
\]
Here and below $i\in\{0,1\}$.
The nonnegative weights $w_{a,i}$ sum to one. If
$\overline\rho=\sum_aq_a\proj{\psi_a}$ and $\overline\rho_B$ is its
partial trace over the first input qubit, then
\begin{equation}\label{eq:partner-marginal}
 \begin{aligned}
 \overline\rho_B&=\sum_{a,i}w_{a,i}\proj{\psi_{a,i}},\\
 \overline\tau:=\sum_{a,i}w_{a,i}\tau_{a,i}
 &=\mathcal N(\overline\rho_B)
 =\Tr_{\C^2}\bigl((\A\otimes\mathcal N)(\overline\rho)\bigr).
 \end{aligned}
\end{equation}
The cross terms disappear in the first identity because
$\langle0|1\rangle=0$, regardless of any overlap between the branch
vectors. The last identity uses trace preservation of $\A$. Thus the
branch ensemble for $\mathcal N$ has exactly the second marginal of
the mean full output.

Write $\sigma_a=(\A\otimes\mathcal N)(\proj{\psi_a})$ and
$\overline\sigma=\sum_aq_a\sigma_a$. The first output marginal has
excited-state population $t(\overline p)=(1-\gamma)\overline p+\gamma\nu$;
here trace preservation of $\mathcal N$ identifies that marginal with
$\A$ applied to the first input marginal. Apply the diagonal-entropy
bound \eqref{eq:diagonal} in the product of the first output's
computational basis and an eigenbasis of $\overline\tau$. The resulting
joint distribution has marginals $(1-t(\overline p),t(\overline p))$
and the eigenvalue distribution of $\overline\tau$. Classical
subadditivity, as proved in \Cref{sec:holevo}, gives
\begin{equation}\label{eq:partner-mean-upper}
 \ent(\overline\sigma)
 \le h(t(\overline p))+\ent(\overline\tau).
\end{equation}
On the other hand, Lemma~\ref{lem:partner-branch} and the convexity in
Lemma~\ref{lem:e-convex} imply
\begin{equation}\label{eq:partner-ensemble-lower}
 \begin{split}
 \sum_aq_a\ent(\sigma_a)
 &\ge\sum_aq_ae(p_a)+\sum_{a,i}w_{a,i}\ent(\tau_{a,i})\\
 &\ge e(\overline p)+\sum_{a,i}w_{a,i}\ent(\tau_{a,i}).
 \end{split}
\end{equation}
Subtracting \eqref{eq:partner-ensemble-lower} from
\eqref{eq:partner-mean-upper} yields
\begin{align*}
 \ent(\overline\sigma)-\sum_aq_a\ent(\sigma_a)
 &\le h(t(\overline p))-e(\overline p)\\
 &\quad+\ent(\overline\tau)-\sum_{a,i}w_{a,i}\ent(\tau_{a,i})\\
 &\le (\ln2)\bigl(\chi(\A)+\chi(\mathcal N)\bigr).
\end{align*}
For the last line, the one-use case of \Cref{thm:main} bounds the scalar
term, and \eqref{eq:holevo} bounds the Holevo information of the branch
ensemble, whose mean is \eqref{eq:partner-marginal}. Taking the
supremum proves the upper bound in \eqref{eq:partner-holevo}.
For the reverse bound, take the phase-pair ensemble attaining
$\chi(\A)$ and, for each $\varepsilon>0$, a finite ensemble with
Holevo information greater than $\chi(\mathcal N)-\varepsilon$.
Their product ensemble has the sum of their Holevo informations,
because both the mean outputs and individual outputs are tensor
products. Taking the supremum proves \eqref{eq:partner-holevo}.

For minimum output entropy, concavity \eqref{eq:concavity} shows that
some pure input attains the minimum of any finite-dimensional channel:
a pure-state decomposition cannot have every component output entropy
larger than that of the mixed input. For $\A$, the pure qubit calculation
\eqref{eq:single-entropy} therefore gives
\begin{equation}\label{eq:minimum-output-value}
 S_{\min}(\A)=\min_{0\le p\le1}e(p)
              =g\bigl(\gamma\nu(1-\nu)\bigr).
\end{equation}
Indeed, the completed square in \eqref{eq:v-main} is minimized at
$p=\nu$, and \Cref{lem:binary-spectrum} implies that $g$ is
nondecreasing. Any pure qubit state with that population attains the
minimum, including at the parameter endpoints.
The right-hand side of \eqref{eq:partner-branch} is at least
$S_{\min}(\A)+S_{\min}(\mathcal N)$. Hence the same is true of every
output entropy of every pure input for $\A\otimes\mathcal N$, and concavity
extends the lower bound to mixed inputs. Tensoring minimizing input
states for the two factors gives the reverse inequality, since entropy
is additive on product states. This proves \eqref{eq:partner-min}.

Finally, a permutation of tensor factors identifies
$(\A\otimes\mathcal N)^{\otimes n}$ with
$\A^{\otimes n}\otimes\mathcal N^{\otimes n}$ up to input and output
unitaries, which do not change $\chi$. Repeated application of
\eqref{eq:partner-holevo} gives, for every integer $n\ge1$,
\begin{equation}\label{eq:partner-blocks}
 \chi\bigl((\A\otimes\mathcal N)^{\otimes n}\bigr)
 =n\chi(\A)+\chi(\mathcal N^{\otimes n}).
\end{equation}
At each step the partner is a tensor product of the remaining copies
of $\A$ with $\mathcal N^{\otimes n}$, hence is again finite dimensional
and completely positive trace preserving. Applying
\eqref{eq:regularization} to \eqref{eq:partner-blocks} gives
\[
 C(\A\otimes\mathcal N)
 =\chi(\A)+\sup_{n\ge1}\frac1n\chi(\mathcal N^{\otimes n})
 =C(\A)+C(\mathcal N),
\]
where the last equality uses $C(\A)=\chi(\A)$ from \Cref{thm:main}.
This proves \eqref{eq:partner-capacity}. In particular, no restriction
has been placed on entanglement among the inputs to the partner's
copies in \eqref{eq:partner-blocks}.
\end{proof}

\section{Ordinary damping, the unital case, and endpoints}
\label{sec:special}

The population convention makes several special cases transparent.
Throughout this section, capacities and Holevo information are in bits,
while $h$ and $g$ remain in nats.

Tang, Zhu, Bai, and Wang
\cite[Corollary~5.1 and Theorem~5.5]{TangEtAl2026} established the
ordinary-damping capacity and its additivity with arbitrary
finite-dimensional partners. The following specialization recovers
their capacity formula in the present parameters.

\begin{corollary}[Ordinary amplitude damping]\label{cor:ordinary}
For $\gamma\in[0,1]$, let $\A_\gamma$ be the qubit channel with Kraus
operators
\[
 K_0=\diag(1,\sqrt{1-\gamma}),\qquad
 K_1=\sqrt\gamma\,\ket0\bra1.
\]
For every integer $n\ge1$,
\[
 \chi(\A_\gamma^{\otimes n})=n C(\A_\gamma)
 =\frac n{\ln2}\max_{0\le p\le1}
 \{h((1-\gamma)p)-g(\gamma(1-\gamma)p^2)\}.
\]
\end{corollary}
\begin{proof}
The Kraus operators give \eqref{eq:channel} with $\nu=0$.
The conclusion follows from \Cref{thm:main} after this substitution.
\end{proof}

The scalar maximum and the opposite-phase attaining ensemble agree
with Giovannetti--Fazio
\cite[Section~III~A, Eqs.~(42)--(43)]{GiovannettiFazio2005}, whose
transmissivity is $1-\gamma$ in our convention. That one-use
optimization becomes the unrestricted classical capacity through the
finite-power identity.

Let $\mathsf X=\bigl(\begin{smallmatrix}0&1\\1&0\end{smallmatrix}\bigr)$.
Direct substitution in \eqref{eq:channel} gives
\[
 \A_{\gamma,1-\nu}(\rho)
 =\mathsf X\A_{\gamma,\nu}(\mathsf X\rho\mathsf X)\mathsf X.
\]
Unitary conjugations preserve the Holevo optimization, so the capacity
is invariant under $\nu\mapsto1-\nu$. In particular, $\nu=1$ has the
ordinary-channel capacity as well, with the signal population replaced
by $1-p$.

The next formula also follows from King's general unital-qubit theorem
\cite{King2002}. We recover it directly from the present scalar
optimization.

\begin{corollary}[Unital midpoint]\label{cor:unital}
For every $\gamma\in[0,1]$,
\[
 C(\A_{\gamma,1/2})
 =1-\frac1{\ln2}g(\gamma/4)
 =1-\frac1{\ln2}h\!\left(\frac{1+\sqrt{1-\gamma}}2\right).
\]
The pair $(\ket0\pm\ket1)/\sqrt2$ attains the one-use Holevo
information, and its independent products attain every finite-power
Holevo optimum.
\end{corollary}
\begin{proof}
Write the one-use objective in nats as
\[
 f(p)=h((1-\gamma)p+\gamma/2)
       -g\bigl(\gamma/4+\gamma(1-\gamma)(p-1/2)^2\bigr).
\]
It is concave: its first term is concave, and its second term is the
negative of the convex function in \Cref{lem:e-convex}. Also
$f(1-p)=f(p)$, using $h(1-t)=h(t)$. Thus
$f(1/2)\ge(f(p)+f(1-p))/2=f(p)$. Evaluating the maximum at $p=1/2$
and using \Cref{thm:main} proves the formulas and the attainment claim.
\end{proof}

At $\gamma=0$, the channel is the identity, $v_{\gamma,\nu}=0$, and
the capacity is one bit per use. At $\gamma=1$, it replaces every
input by $\diag(1-\nu,\nu)$, so its capacity is zero: the two terms
in the objective both equal $h(\nu)$. These conclusions include all
four corners of the parameter square.

\appendix
\section{Alternative Hessian proofs}
\label{sec:alternative}

The main proof obtains scalar concavity and matrix log-determinant
convexity from the same contraction argument. Here we give direct
Hessian proofs of the two analytic statements. They provide an
alternative route to \Cref{lem:scalar,lem:logdet}; the entropy integral,
Schur complement, sign averaging, and rectangular embedding in
\Cref{sec:block} are unchanged.

\subsection{The Hessian of the scalar deficit}

Recall the function $F$ from \eqref{eq:block-F}. Its homogeneity follows
by the scaling cancellation already proved in \Cref{lem:scalar}.
We verify concavity directly on the probability simplex, where
$F(x,y,z)=H(x,y,z)-g(yz)$, using the derivatives
\eqref{eq:g-derivatives}. For $0<u<1/4$, write
$t=\sqrt{1-4u}$ and $L=\atanh(t)/t$.

\begin{proof}[Alternative proof of \Cref{lem:scalar}]
In the simplex interior, set $x=1-y-z$ and $u=yz$.
The negative Hessian of $F(1-y-z,y,z)$ is
\[
 J=\begin{pmatrix}
 \dfrac1y+\dfrac1x+z^2g''(u)&\dfrac1x+g'(u)+ug''(u)\\[5pt]
 \dfrac1x+g'(u)+ug''(u)&\dfrac1z+\dfrac1x+y^2g''(u)
 \end{pmatrix}.
\]
Expanding its determinant and substituting \eqref{eq:g-derivatives}
and $x+y+z=1$ yields
\begin{align}
 \det J
 &=\frac1{xu}+\frac{y+z-4u}{x}g''(u)
   -\frac2xg'(u)-g'(u)^2-2ug'(u)g''(u)\notag\\
 &=\frac{1-4uL^2}{ut^2}>0.
 \label{eq:block-scalar-determinant}
\end{align}
Indeed, for $0<t<1$,
\[
 \atanh t<\frac{t}{\sqrt{1-t^2}},
\]
as both sides vanish at zero and their derivatives are respectively
$(1-t^2)^{-1}$ and $(1-t^2)^{-3/2}$. Since $4u=1-t^2$, this proves
the strict inequality in \eqref{eq:block-scalar-determinant}.
Furthermore,
\[
 \begin{pmatrix}y&z\end{pmatrix}J\begin{pmatrix}y\\z\end{pmatrix}
 =\frac{y+z}{x}
   +4u\left(ug''(u)+\frac12g'(u)\right)>0.
\]
A real symmetric $2\times2$ matrix with positive determinant has
eigenvalues of the same sign; this positive quadratic form forces
both to be positive. Hence $F$ is concave in the simplex interior,
and continuity extends concavity to the closed simplex.

To prove \eqref{eq:block-superadditive}, omit zero triples and put
$m_i=x_i+y_i+z_i$, $m=\sum_i m_i$. For $m>0$, homogeneity and
simplex concavity give
\[
 \sum_i F(x_i,y_i,z_i)
 =m\sum_i\frac{m_i}{m}
      F\!\left(\frac{x_i}{m_i},\frac{y_i}{m_i},\frac{z_i}{m_i}\right)
 \le F\!\left(\sum_i x_i,\sum_i y_i,\sum_i z_i\right).
\]
If $m=0$, all terms vanish.

\end{proof}

\subsection{The Hessian of the log-determinant difference}

The contraction proof avoids choosing singular-vector coordinates.
The following calculation instead diagonalizes the fixed matrix and
expresses the second derivative as a sum of nonnegative squares.
It proves \Cref{lem:logdet} directly for the same domain, including
singular $Z$.

\begin{proof}[Alternative proof of \Cref{lem:logdet}]
At any point $(D_0,E_0)$, use the fixed congruences
\[
 D=D_0^{1/2}\widetilde D D_0^{1/2},\qquad
 E=E_0^{1/2}\widetilde E E_0^{1/2},\qquad
 W=D_0^{-1/2}ZE_0^{-1/2}.
\]
The determinant factors cancel to give
$G_Z(D,E)=G_W(\widetilde D,\widetilde E)$.
These invertible linear changes preserve affine lines. Independent
unitary changes in the two spaces then put $W$ in singular-value form
$Y=\diag(v_1,\ldots,v_N)$ with $v_i\ge0$. It therefore suffices to
prove nonnegativity of the second derivative at $(I,I)$ for this $Y$
and arbitrary Hermitian directions $P,Q$.

Set $R=(I+YY^*)^{-1}$ and $N_0=P-YQY^*$.
Using $(I+sQ)^{-1}=I-sQ+s^2Q^2+O(s^3)$ gives
\begin{equation}\label{eq:block-G-hessian}
 \left.\frac{\mathrm d^2}{\mathrm ds^2}
 G_Y(I+sP,I+sQ)\right|_{s=0}
 =\Tr\bigl(P^2+2RYQ^2Y^*-RN_0RN_0\bigr).
\end{equation}
For completeness, the log-determinant expansion used here is
\[
 \ln\det(H+sK+s^2L_0)
 =\ln\det H+s\Tr(H^{-1}K)
 +s^2\left(\Tr(H^{-1}L_0)
       -\frac12\Tr(H^{-1}KH^{-1}K)\right)+O(s^3),
\]
for $H>0$ and Hermitian $K,L_0$. It follows by factoring $H^{1/2}$
and applying the scalar logarithm expansion to the eigenvalues of the
resulting Hermitian increment. No commutativity is assumed.

Write $r_i=(1+v_i^2)^{-1}$ and $c_i=v_i/(1+v_i^2)$.
Expanding the trace in \eqref{eq:block-G-hessian} in entries gives
\begin{align*}
 \sum_{i,j}\bigl[&(1-r_ir_j)|P_{ij}|^2
 +\{2(1-r_i)-(1-r_i)(1-r_j)\}|Q_{ij}|^2\\
 &\hspace{34mm}+2c_ic_j\operatorname{Re}
       (P_{ij}\overline{Q_{ij}})\bigr].
\end{align*}
Since $Q$ is Hermitian, $|Q_{ij}|^2=|Q_{ji}|^2$.
Pairing the $(i,j)$ and $(j,i)$ terms replaces their coefficient by
\[
 \frac12\bigl[2(1-r_i)+2(1-r_j)
              -2(1-r_i)(1-r_j)\bigr]=1-r_ir_j.
\]
This identity also holds when $i=j$. Completing the square therefore
rewrites the second derivative as
\begin{equation}\label{eq:block-G-squares}
 \sum_{i,j}\left[
 (1-r_ir_j-c_ic_j)(|P_{ij}|^2+|Q_{ij}|^2)
       +c_ic_j|P_{ij}+Q_{ij}|^2\right].
\end{equation}
Every coefficient is nonnegative: $c_ic_j\ge0$, and
\[
 1-r_ir_j-c_ic_j
 =\frac{(v_i-v_j)^2+v_iv_j+v_i^2v_j^2}
        {(1+v_i^2)(1+v_j^2)}\ge0.
\]
Thus the second derivative is nonnegative in every Hermitian
direction at every positive definite pair. The domain is convex,
so the asserted joint convexity follows along each line segment.
\end{proof}

\begin{thebibliography}{99}

\bibitem{Audenaert2008}
K.~M.~R. Audenaert,
``On a norm compression inequality for $2\times N$ partitioned block matrices,''
\emph{Linear Algebra and its Applications} \textbf{428}, 781--795 (2008).
doi: \doi{10.1016/j.laa.2007.08.007}.

\bibitem{BennettEtAl2002}
C.~H. Bennett, P.~W. Shor, J.~A. Smolin, and A.~V. Thapliyal,
``Entanglement-assisted capacity of a quantum channel and the reverse Shannon theorem,''
\emph{IEEE Transactions on Information Theory} \textbf{48}(10),
2637--2655 (2002).
doi: \doi{10.1109/TIT.2002.802612}.
Preprint: arXiv:\href{https://arxiv.org/abs/quant-ph/0106052v2}{quant-ph/0106052v2}.

\bibitem{Berry2005}
D.~W. Berry,
``Qubit channels that achieve capacity with two states,''
\emph{Physical Review A} \textbf{71}, 032334 (2005).
doi: \doi{10.1103/PhysRevA.71.032334}.

\bibitem{Cortese2002}
J. Cortese,
\emph{Relative entropy and single qubit Holevo--Schumacher--Westmoreland channel capacity},
arXiv:\href{https://arxiv.org/abs/quant-ph/0207128}{quant-ph/0207128} (2002).

\bibitem{Fang2026}
K. Fang,
\emph{A pretty-tight converse on the classical capacity of generalized amplitude-damping channels},
arXiv:\href{https://arxiv.org/abs/2609.14621v1}{2609.14621v1} (2026).

\bibitem{GiovannettiFazio2005}
V. Giovannetti and R. Fazio,
``Information-capacity description of spin-chain correlations,''
\emph{Physical Review A} \textbf{71}, 032314 (2005).
doi: \doi{10.1103/PhysRevA.71.032314}.

\bibitem{Hastings2009}
M.~B. Hastings,
``Superadditivity of communication capacity using entangled inputs,''
\emph{Nature Physics} \textbf{5}, 255--257 (2009).
doi: \doi{10.1038/nphys1224}.

\bibitem{HausladenEtAl1996}
P. Hausladen, R. Jozsa, B. Schumacher, M. Westmoreland, and W.~K. Wootters,
``Classical information capacity of a quantum channel,''
\emph{Physical Review A} \textbf{54}, 1869--1876 (1996).
doi: \doi{10.1103/PhysRevA.54.1869}.

\bibitem{Hiai2016}
F. Hiai,
``Concavity of certain matrix trace and norm functions.~II,''
\emph{Linear Algebra and its Applications} \textbf{496}, 193--220 (2016).
doi: \doi{10.1016/j.laa.2015.12.032}.
Preprint: arXiv:\href{https://arxiv.org/abs/1507.00853v3}{1507.00853v3}.

\bibitem{Holevo1998}
A.~S. Holevo,
``The capacity of the quantum channel with general signal states,''
\emph{IEEE Transactions on Information Theory} \textbf{44}(1),
269--273 (1998).
doi: \doi{10.1109/18.651037}.
Preprint: arXiv:\href{https://arxiv.org/abs/quant-ph/9611023v1}{quant-ph/9611023v1}.

\bibitem{HouFang2007}
L.-Z. Hou and M.-F. Fang,
``The Holevo capacity of a generalized amplitude-damping channel,''
\emph{Chinese Physics} \textbf{16}(7), 1843--1847 (2007).
doi: \doi{10.1088/1009-1963/16/7/006}.

\bibitem{KhatriEtAl2020}
S. Khatri, K. Sharma, and M.~M. Wilde,
``Information-theoretic aspects of the generalized amplitude-damping channel,''
\emph{Physical Review A} \textbf{102}, 012401 (2020).
doi: \doi{10.1103/PhysRevA.102.012401}.

\bibitem{King2002}
C. King,
``Additivity for unital qubit channels,''
\emph{Journal of Mathematical Physics} \textbf{43}, 4641--4653 (2002).
doi: \doi{10.1063/1.1500791}.

\bibitem{King2003}
C. King,
``Inequalities for trace norms of $2\times2$ block matrices,''
\emph{Communications in Mathematical Physics} \textbf{242}, 531--545 (2003).
doi: \doi{10.1007/s00220-003-0955-9}.

\bibitem{LeditzkyEtAl2018}
F. Leditzky, E. Kaur, N. Datta, and M.~M. Wilde,
``Approaches for approximate additivity of the Holevo information of quantum channels,''
\emph{Physical Review A} \textbf{97}, 012332 (2018).
doi: \doi{10.1103/PhysRevA.97.012332}.

\bibitem{SchumacherWestmoreland1997}
B. Schumacher and M.~D. Westmoreland,
``Sending classical information via noisy quantum channels,''
\emph{Physical Review A} \textbf{56}(1), 131--138 (1997).
doi: \doi{10.1103/PhysRevA.56.131}.

\bibitem{Shor2002}
P.~W. Shor,
``Additivity of the classical capacity of entanglement-breaking quantum channels,''
\emph{Journal of Mathematical Physics} \textbf{43}, 4334--4340 (2002).
doi: \doi{10.1063/1.1498000}.

\bibitem{TangEtAl2026}
Z. Tang, C. Zhu, G. Bai, and X. Wang,
\emph{Classical capacity and entanglement cost of the amplitude damping channel},
arXiv:\href{https://arxiv.org/abs/2609.28592v1}{2609.28592v1} (2026).

\bibitem{Wootters1998}
W.~K. Wootters,
``Entanglement of formation of an arbitrary state of two qubits,''
\emph{Physical Review Letters} \textbf{80}, 2245--2248 (1998).
doi: \doi{10.1103/PhysRevLett.80.2245}.

\end{thebibliography}

\end{document}
